% ============================================================
%  article.tex  —  TEMPLATE for a new journal article
%
%  To add a new article to the journal:
%    1.  Copy the entire  template/article/  directory into
%        the project root, renaming it (e.g. my-new-article/).
%    2.  Edit this file: fill in the \articleheader fields and
%        write your content (or \input section files).
%    3.  In main.tex, add:  % ============================================================
%  article.tex  —  TEMPLATE for a new journal article
%
%  To add a new article to the journal:
%    1.  Copy the entire  template/article/  directory into
%        the project root, renaming it (e.g. my-new-article/).
%    2.  Edit this file: fill in the \articleheader fields and
%        write your content (or \input section files).
%    3.  In main.tex, add:  % ============================================================
%  article.tex  —  TEMPLATE for a new journal article
%
%  To add a new article to the journal:
%    1.  Copy the entire  template/article/  directory into
%        the project root, renaming it (e.g. my-new-article/).
%    2.  Edit this file: fill in the \articleheader fields and
%        write your content (or \input section files).
%    3.  In main.tex, add:  % ============================================================
%  article.tex  —  TEMPLATE for a new journal article
%
%  To add a new article to the journal:
%    1.  Copy the entire  template/article/  directory into
%        the project root, renaming it (e.g. my-new-article/).
%    2.  Edit this file: fill in the \articleheader fields and
%        write your content (or \input section files).
%    3.  In main.tex, add:  \input{my-new-article/article}
%    4.  Add a \contentsentry in the Table‑of‑Contents block.
% ============================================================

\label{art:CHANGEME}

\articleheader
  {TITLE OF YOUR ARTICLE}
  {YOUR NAME}
  {Constructor University Bremen}
  {ABSTRACT — A concise summary of the article's contributions,
   methodology, and key results.}

% ── Optional: Acknowledgements ────────────────────────────
% \vspace{0.3cm}
% \noindent{\small\textbf{\color{cujblue}Acknowledgements.}\ %
% Thank collaborators, advisors, funding sources, etc.}
% \vspace{0.6cm}

% ── Article body ──────────────────────────────────────────

\section{Introduction}

Write your introduction here.

% You can split content into separate files:
% \input{my-new-article/sections/content}

\section{Preliminaries}

Definitions, notation, background material.

\section{Main Results}

Core content: theorems, proofs, examples, figures.

\section{Conclusion}

Summary and future work.

% ── References ────────────────────────────────────────────
\articlereferences

\begin{enumerate}[label={[\arabic*]}, leftmargin=2em, itemsep=0.25em]
  \item A.~Author, ``Title,'' \textit{Journal}, vol.~X, pp.~1--10, 2025.
  \item B.~Author, \textit{Book Title}, Publisher, 2025.
\end{enumerate}

%    4.  Add a \contentsentry in the Table‑of‑Contents block.
% ============================================================

\label{art:CHANGEME}

\articleheader
  {TITLE OF YOUR ARTICLE}
  {YOUR NAME}
  {Constructor University Bremen}
  {ABSTRACT — A concise summary of the article's contributions,
   methodology, and key results.}

% ── Optional: Acknowledgements ────────────────────────────
% \vspace{0.3cm}
% \noindent{\small\textbf{\color{cujblue}Acknowledgements.}\ %
% Thank collaborators, advisors, funding sources, etc.}
% \vspace{0.6cm}

% ── Article body ──────────────────────────────────────────

\section{Introduction}

Write your introduction here.

% You can split content into separate files:
% % ============================================================
%  content.tex  —  Template section file
%  Replace this with your article's section content.
% ============================================================

\section{Your Section Title}

Write your content here.  You can use all standard \LaTeX\ commands
as well as the journal environments:

\begin{definitionbox}
  State your definition here.
\end{definitionbox}

\begin{theorembox}
  State your theorem here.
\end{theorembox}

\begin{notebox}
  Add a side note or remark here.
\end{notebox}


\section{Preliminaries}

Definitions, notation, background material.

\section{Main Results}

Core content: theorems, proofs, examples, figures.

\section{Conclusion}

Summary and future work.

% ── References ────────────────────────────────────────────
\articlereferences

\begin{enumerate}[label={[\arabic*]}, leftmargin=2em, itemsep=0.25em]
  \item A.~Author, ``Title,'' \textit{Journal}, vol.~X, pp.~1--10, 2025.
  \item B.~Author, \textit{Book Title}, Publisher, 2025.
\end{enumerate}

%    4.  Add a \contentsentry in the Table‑of‑Contents block.
% ============================================================

\label{art:CHANGEME}

\articleheader
  {TITLE OF YOUR ARTICLE}
  {YOUR NAME}
  {Constructor University Bremen}
  {ABSTRACT — A concise summary of the article's contributions,
   methodology, and key results.}

% ── Optional: Acknowledgements ────────────────────────────
% \vspace{0.3cm}
% \noindent{\small\textbf{\color{cujblue}Acknowledgements.}\ %
% Thank collaborators, advisors, funding sources, etc.}
% \vspace{0.6cm}

% ── Article body ──────────────────────────────────────────

\section{Introduction}

Write your introduction here.

% You can split content into separate files:
% % ============================================================
%  content.tex  —  Template section file
%  Replace this with your article's section content.
% ============================================================

\section{Your Section Title}

Write your content here.  You can use all standard \LaTeX\ commands
as well as the journal environments:

\begin{definitionbox}
  State your definition here.
\end{definitionbox}

\begin{theorembox}
  State your theorem here.
\end{theorembox}

\begin{notebox}
  Add a side note or remark here.
\end{notebox}


\section{Preliminaries}

Definitions, notation, background material.

\section{Main Results}

Core content: theorems, proofs, examples, figures.

\section{Conclusion}

Summary and future work.

% ── References ────────────────────────────────────────────
\articlereferences

\begin{enumerate}[label={[\arabic*]}, leftmargin=2em, itemsep=0.25em]
  \item A.~Author, ``Title,'' \textit{Journal}, vol.~X, pp.~1--10, 2025.
  \item B.~Author, \textit{Book Title}, Publisher, 2025.
\end{enumerate}

%    4.  Add a \contentsentry in the Table‑of‑Contents block.
% ============================================================

\label{art:CHANGEME}

\articleheader
  {TITLE OF YOUR ARTICLE}
  {YOUR NAME}
  {Constructor University Bremen}
  {ABSTRACT — A concise summary of the article's contributions,
   methodology, and key results.}

% ── Optional: Acknowledgements ────────────────────────────
% \vspace{0.3cm}
% \noindent{\small\textbf{\color{cujblue}Acknowledgements.}\ %
% Thank collaborators, advisors, funding sources, etc.}
% \vspace{0.6cm}

% ── Article body ──────────────────────────────────────────

\section{Introduction}

Write your introduction here.

% You can split content into separate files:
% % ============================================================
%  content.tex  —  Template section file
%  Replace this with your article's section content.
% ============================================================

\section{Your Section Title}

Write your content here.  You can use all standard \LaTeX\ commands
as well as the journal environments:

\begin{definitionbox}
  State your definition here.
\end{definitionbox}

\begin{theorembox}
  State your theorem here.
\end{theorembox}

\begin{notebox}
  Add a side note or remark here.
\end{notebox}


\section{Preliminaries}

Definitions, notation, background material.

\section{Main Results}

Core content: theorems, proofs, examples, figures.

\section{Conclusion}

Summary and future work.

% ── References ────────────────────────────────────────────
\articlereferences

\begin{enumerate}[label={[\arabic*]}, leftmargin=2em, itemsep=0.25em]
  \item A.~Author, ``Title,'' \textit{Journal}, vol.~X, pp.~1--10, 2025.
  \item B.~Author, \textit{Book Title}, Publisher, 2025.
\end{enumerate}
